\section{Methods}
\label{sec:methods}

% Estimator zoo + attribution protocol + block ladder construction.
% Nothing here reports an evaluation loss; every number in this section is
% either a definition or a design-justifying measurement on a non-scored
% panel, disclosed as such.

The methods divide into the estimators and the constructions.
Section~\ref{sec:estimators} defines the estimator set --- one
unpenalized member, the penalized linear family, and two tree ensembles
--- together with the walk-forward protocol and the causal tuning rule
every tuned arm obeys. Section~\ref{sec:attribution} defines the
attribution protocol. Sections~\ref{sec:dense_weak_twoblock} through
\ref{sec:three_block} build the block ladder the paper ends on: a
two-block ridge matched to the measured coefficient structure, an
explicit product block distilled from the tree fits, and the three-block
ridge that carries both. Every estimator is fit on the shared panel of
Section~\ref{sec:data_prep}, scored under the contract of
Section~\ref{sec:evaluation}, and subject to the hyperparameter
discipline of Section~\ref{sec:dp_benchmark}: no quantity is selected on
evaluation loss.

\subsection{Estimators and Causal Tuning}
\label{sec:estimators}

\paragraph{Minimum-norm least squares.} The unpenalized estimator is
\emph{minimum-norm least squares}: among all coefficient vectors that
minimize the training-window sum of squares, the unique one of minimal
$\ell_2$ norm, computed through the pseudoinverse of the window gram at
the standard machine-precision cutoff. This is the $\alpha \to 0^{+}$
limit of ridge regression, it is defined at any design rank, and it
carries no penalty and no tuned quantity of any kind --- the strongest
attribution claim an unpenalized estimator supports. The rank cutoff is
not a tuned quantity either: the centered gram's spectrum is gapped by
more than five orders of magnitude around the retention threshold (on
the widest design, largest discarded eigenvalue $1.8 \times 10^{-11}$
relative, weakest retained $6.2 \times 10^{-6}$, with no eigenvalue
between them at any bar), so every tolerance across four decades retains
the identical subspace and reproduces the reported forecasts bitwise. Each arm is refit at
every bar over a rolling window of $24{,}000$ bars ($500$ trading days).
On a full-rank design (the benchmark's $52$ columns) minimum-norm least
squares coincides exactly with ordinary least squares. An earlier pass of
this campaign used the plain normal-equations solve, whose solution is
arbitrary in near-null directions of the design; two single-window
incidents documented in Section~\ref{sec:buckets} trace to that path, and
its results are archived. Because the pseudoinverse is defined at any
rank, the estimator requires no rank guards; the only design
intervention retained is the go-live participation rule of
Section~\ref{sec:dp_prep}, which is a statement about data existence,
not numerics.

\paragraph{Estimators.}
The design matrix $X$ has one row per bar in the current training window
and one column per feature: the HAR ladder of the target, the calendar
features, and the full exogenous block of $523$ columns. Let $y$
be the vector of targets over the same window, and let $x_s$ denote row $s$
of $X$. The two estimators solve, at each refit,
\begin{align*}
\hat\beta^{\mathrm{ridge}} &= \arg\min_{\beta}
  \sum_{s} \left( y_s - x_s^\top \beta \right)^2
  + \alpha \,\lVert \beta \rVert_2^2, \\
\hat\beta^{\mathrm{lasso}} &= \arg\min_{\beta}
  \sum_{s} \left( y_s - x_s^\top \beta \right)^2
  + \lambda \,\lVert \beta \rVert_1,
\end{align*}
where the sums run over the bars of the training window and the penalties
$\alpha$ and $\lambda$ are held constant between the periodic re-selections
described below. The $\ell_1$ penalty produces a coefficient vector
supported on a subset of columns; the $\ell_2$ penalty produces a
coefficient vector with all entries nonzero.

\paragraph{Walk-forward protocol.}
The panel is computed as 100 walk-forward chunks over a rolling
500-trading-day ($24{,}000$-bar) training window. Every arm is refit
\emph{every bar}: at each bar the oldest bar leaves the window, the newest
enters, and the coefficients are re-estimated before the next forecast is
issued. Each arm re-selects its own penalty periodically and causally: every
250 solves, the current training window is split forward into a fit block, a
25-bar embargo (covering the longest moving average entering the features),
and a 125-bar validation tail, and the penalty value minimizing validation
MSE on that tail is adopted. No forecast's hyperparameters ever see its own
future. The lasso arm is solved as an exact every-bar online homotopy: the
solution and active set at each bar are obtained by exactly updating the
previous bar's solution for the one-row change in the window, not by a
warm-started batch solver terminated at a tolerance, so the head-to-head is
not confounded with optimization error. The ridge arm runs on an exact
rank-one update path.
% Estimator engineering (rank-one algebra, homotopy, BlockRidge solver) is
% appendix material per the restructure plan.

\paragraph{Tree ensembles.} The nonlinear members of the estimator set
are the gradient-boosted tree ensembles LightGBM \citep{Ke2017} and
XGBoost \citep{Chen2016}. The comparison battery fully crosses the
estimators of this subsection with the nine feature buckets on the
shared panel ($n = 218{,}934$ thirty-minute bars, identical rows in
every cell). Every arm, tree and linear alike, reselects its own
hyperparameters on the same periodic causal schedule: each reselection
uses only data available before the reselection date, so no forecast's
configuration uses information from that forecast's future. The shared
comparator is the per-bar OLS--HAR incumbent
(Section~\ref{sec:dp_benchmark}).


\subsection{The Attribution Protocol}
\label{sec:attribution}

\paragraph{Protocol.} The attribution is computed as follows.
\begin{enumerate}
\item \emph{Arms.} The incumbent arm's design matrix contains the HAR ladder
columns and the calendar dummies. Each bucket arm's design matrix contains
those backbone columns plus the bucket's own columns, expanded through the
same moving-average ladder and availability indicators as the joint arm. The
joint arm's design matrix contains the backbone columns plus all $41$
exogenous series ($\approx 526$ columns after the moving-average ladder and
availability indicators).
\item \emph{Fit.} Each arm is fit by the rolling estimator of
Section~
ef{sec:estimators} and produces
one out-of-sample forecast per bar.
\item \emph{Loss.} Each arm is scored by per-bar QLIKE on the shared panel
($n = 218{,}934$ thirty-minute bars, identical rows in every cell). The
reported QLIKE is the mean of the arm's per-bar losses.
\item \emph{Increments.} $\Delta$ is the arm's QLIKE minus the incumbent's
QLIKE. $\Delta_{\text{own}}$ is the arm's QLIKE minus the same estimator's
own no-exogenous baseline ($0.13445$).
\item \emph{DM test.} For each arm, the per-bar loss differential against the
incumbent is $d_t = \ell^{\text{arm}}_t - \ell^{\text{inc}}_t$. The
Diebold--Mariano statistic \citep{DieboldMariano1995} is the $t$-ratio of
$\bar{d}$ to its Newey--West HAC standard error, with the
Harvey--Leybourne--Newbold small-sample correction
\citep{HarveyLeybourneNewbold1997}. Negative $t$ favours the arm.
\item \emph{Overlap.} The overlap ratio is
\[
R \;=\; \frac{\sum_{b} \Delta_{\text{own}}(b)}{\Delta_{\text{own}}(\text{joint})},
\]
the sum of the single-bucket within-class increments divided by the joint
arm's within-class increment, both measured against the same estimator's
no-exogenous baseline.
\end{enumerate}

\subsection{The Two-Block Ridge}
\label{sec:dense_weak_twoblock}

The head-to-head leaves two facts that a single ridge penalty cannot
satisfy at once. First, penalizing the HAR-plus-calendar design alone
raises QLIKE relative to plain OLS (the battery's no-exogenous control arm:
ridge $t = +5.0$ against the incumbent): the HAR columns are few,
well-conditioned, and individually strong, so any shrinkage on them is pure
bias. Second, the wide exogenous design requires heavy shrinkage --- that is
the head-to-head's result. A scalar penalty applies the same $\alpha$ to both
blocks:
a value large enough for the hundreds of weak exogenous columns over-shrinks
the HAR columns, and a value small enough for the HAR columns under-shrinks
the exogenous columns. The resolution is a penalty that differs by block.

The section's estimator is a \textbf{two-block ridge}: one jointly solved
least-squares problem with a block-diagonal Tikhonov penalty,
\[
(\hat\beta_1, \hat\beta_2)
= \arg\min_{\beta_1, \beta_2}
  \left\lVert y - X_1 \beta_1 - X_2 \beta_2 \right\rVert_2^2
  + \alpha_1 \lVert \beta_1 \rVert_2^2
  + \alpha_2 \lVert \beta_2 \rVert_2^2,
\qquad (\alpha_1, \alpha_2) = (1, 100).
\]
Block 1 ($X_1$) is the HAR ladder of the target, at $\alpha_1 = 1$. Block 2
($X_2$) is the same HAR averaging construction applied to \emph{every}
exogenous feature --- HAR-of-\texttt{all\_features}, the full wide basis ---
at $\alpha_2 = 100$. The two blocks are not fit separately and stacked; the
coefficients are estimated in one solve. The per-block penalties reduce to a
single ridge by column scaling: with
$\tilde X = \bigl[\, X_1 / \sqrt{\alpha_1} \;\; X_2 / \sqrt{\alpha_2} \,\bigr]$
and $\gamma = (\sqrt{\alpha_1}\, \beta_1,\; \sqrt{\alpha_2}\, \beta_2)$,
\[
\left\lVert y - X_1 \beta_1 - X_2 \beta_2 \right\rVert_2^2
+ \alpha_1 \lVert \beta_1 \rVert_2^2
+ \alpha_2 \lVert \beta_2 \rVert_2^2
= \left\lVert y - \tilde X \gamma \right\rVert_2^2
+ \lVert \gamma \rVert_2^2 ,
\]
so the two-block problem is an ordinary ridge with unit penalty on the
rescaled design, and $\hat\beta_j = \hat\gamma_j / \sqrt{\alpha_j}$ for
$j = 1, 2$.
\pending{how the $(\alpha_1, \alpha_2) = (1, 100)$ pair was fixed ---
grid/validation provenance, and its sensitivity}

The estimator matches the measured coefficient structure clause by clause.
The signal is spread across many columns, so block 2 keeps all of its
columns: nothing is selected. No individual exogenous coefficient is large,
so block 2 is shrunk two orders of magnitude harder than block 1, and the
fit sums many attenuated contributions. The HAR columns are few and strong,
so block 1 is penalized only at $\alpha_1 = 1$; setting $\beta_2 = 0$
reduces the model to ridge at $\alpha = 1$ on the backbone, which differs
from the OLS--HAR incumbent of Section~\ref{sec:data_prep} only by that
nominal penalty. The estimator contains no selection step and no tuned
dial: the penalties are fixed constants, not optimized quantities.

\subsection{The Product Block}
\label{sec:nonlinear_product}
\label{sec:product_block}

The construction enumerates candidate products, selects a subset, and
appends the subset to the linear design as ordinary columns, read by the
same rolling ridge machinery as every other block. The interaction study
that developed the block settled four design choices: which products,
reselected how often, shrunk how hard, and refit how often. Its numbers
were computed on the interaction-study panel --- $218{,}934$ out-of-sample
bars, 2007--2024, a $250$-day rolling window, HAR reference QLIKE
$0.13275$ --- which shares its rows with the battery panel but not its
window or backbone convention. The levels below are therefore not
comparable to the battery tables and are quoted only for the design
decisions they settle.
\pending{confirmation that the paper's production product block is exactly
this construction (candidate space, selection window, column scaling) once
the block is re-run under the paper's panel convention}

\paragraph{Construction.} The block is built in four steps. The constants
below are those frozen in the campaign executor,
\texttt{src/unification.py}, the code that runs the campaign's product
arms; where the interaction study's description and the executor differ,
the executor's construction is the one stated here, because it is what
ran.
\begin{enumerate}
    \item \emph{Candidate set.} The base set is the target's HAR ladder,
    the $41$ exogenous value channels at trailing windows of $1$, $32$,
    and $512$ bars, and the four session-clock columns (the open, close,
    and overnight indicators and the hour of day) --- $133$ columns on
    the interaction-study panel. The candidate set is every pairwise
    product of two base columns, including squares: $8{,}911$ candidate
    products.
    \item \emph{Selection target.} A backbone-only ridge --- the
    HAR-plus-calendar block at penalty $\alpha = 1$, refit every bar ---
    is walked forward over the first $2W$ rows of the panel, where $W$ is
    the arm's rolling fit window: $24{,}000$ bars ($500$ days) under the
    stated convention of Section~\ref{sec:three_block_results}, $12{,}000$ bars
    ($250$ days) under the documented convention. The selection target is
    that ridge's out-of-sample residual on rows $[W, 2W)$.
    \item \emph{Scoring and selection.} Each base column is centered over
    rows $[W, 2W)$. Each candidate is scored by the absolute correlation,
    over those rows, between the product of its two centered parents and
    the selection target. The $k = 100$ candidates with the largest
    scores are kept. Selection runs once; the set is never reselected;
    and every arm carrying the block is evaluated only from row $2W$
    onward, so no scored bar overlaps the selection window.
    \item \emph{Column scaling.} Each kept product is formed bar by bar
    as the product of its two parent columns as they stand in the panel
    (uncentered, carrying the causal prescaling of
    Section~\ref{sec:data_prep}). The product column is then divided by
    its own trailing rolling standard deviation over a $W$-bar window,
    lagged one bar and requiring at least $1{,}000$ observations, with
    the scale floored from below at $0.1$ times the per-bar cross-column
    median of the product columns' rolling standard deviations --- the
    floor bounds the inflation a low-dispersion trailing window can
    produce, the failure mode the scaling discipline of
    Section~\ref{sec:data_prep} guards against. Rows before the first
    valid scale, all inside the first window and before any scored bar,
    reuse the first valid value. A product with zero dispersion over its
    entire trailing window is set to zero on those bars. Every product
    value at bar $t$ is a function of data available at bar $t$.
\end{enumerate}

\paragraph{Set size.} The interaction study's set-size ladder is
unimodal: accuracy peaks at $k = 25$--$100$ of $8{,}911$ and degrades by
$k = 400$. The linear channel of Section~\ref{sec:buckets_tuned} improved
monotonically in the number of columns; the interaction channel peaks and
degrades. The block uses $k = 100$.

\paragraph{Reselection cadence.} Two arms were compared with identical
machinery and identical coefficient refits, differing only in when
selection runs. The static arm selects the $100$ products once, on the
first selection window (February 2001 through October 2003 in calendar
time), and keeps them for the rest of the sample. The dynamic
arm reselects the product set every month. The dynamic arm loses: at
$k = 100$ it trails the static arm at DM $t = -3.21$ under the study's
original scaling, and the sign is negative at every healthy specification
tested. The dynamic set replaces $24$--$31\%$ of its members per month. The
product set is selected once, eighteen years before the end of the sample,
and is never reselected.

\paragraph{Block penalty.} Products of centered features have variance on
the order of the product of their parents' variances, so the product block
and the base block have columns on incommensurate scales.
Under a single shared penalty the product block's measured contribution is
unstable and can be negative; an early $+43\%$ headline for the block was
an artifact of a design clip that was standing in for the missing
block-wise regularization. Under a separate penalty, with the product block
shrunk an order of magnitude harder than the linear block, the hundred
frozen products add $\Delta R^2 = +0.0067$ over the full $516$-column
linear base --- a $+35\%$ relative improvement in residual-space accuracy,
DM $t = +2.71$. Lighter product penalties eliminate the gain.

\paragraph{Coefficient refit cadence.} The comparison holds the frozen
block and the solver fixed and varies only how often the coefficients are
re-estimated. At a monthly coefficient refit, the block's QLIKE increment
is indistinguishable from zero (DM $t = +0.52$). At a daily refit, the
identical frozen block's residual-space increment is $+0.0108$, and the
raw-scale reconstruction gives $\Delta\text{QLIKE} = -0.00093$, DM
$t = +3.56$. The selected products are predominantly fast pairs built from
one-bar quantities. The identity of the products is frozen; their
coefficients are refit frequently.

Assembled, the product block is one hundred columns with four measured
properties --- sparse, frozen, heavily shrunk, frequently refit --- and it
converts the specification statement of Section~\ref{sec:tree_edge} into a
feature block the paper's estimator can carry.

\subsection{The Three-Block Ridge}
\label{sec:three_block}

The final model appends the product block of
Section~\ref{sec:product_block} to the two-block ridge of
Section~\ref{sec:dense_weak_twoblock} as a third block. Each block carries its own ridge penalty inside one rolling solver:
\begin{equation}
    \hat{y}_t \;=\;
    \underbrace{X^{\mathrm{HAR}}_t \beta_1}_{\alpha_1 = 1}
    \;+\;
    \underbrace{X^{\mathrm{exog}}_t \beta_2}_{\alpha_2 = 100}
    \;+\;
    \underbrace{X^{\mathrm{prod}}_t \beta_3}_{\alpha_3 = 1000},
    \label{eq:three_block}
\end{equation}
where block one is the target's own HAR ladder at $\alpha_1 = 1$; block two
is the linear exogenous block of Section~\ref{sec:dense_weak_twoblock} at
$\alpha_2 = 100$; and block three is the frozen product block at
$\alpha_3 = 1000$. The estimator is implemented by column scaling. Each
column in block $b$ is multiplied by $\sqrt{\alpha_1/\alpha_b}$, and a
single ridge with penalty $\alpha_1$ is solved on the scaled design; a
coefficient $\tilde{\beta}$ on a scaled column corresponds to
$\beta = \sqrt{\alpha_1/\alpha_b}\,\tilde{\beta}$ on the raw column, and
the shared penalty $\alpha_1 \|\tilde{\beta}\|^2$ on the scaled
coefficients equals the per-block penalties $\alpha_b \|\beta\|^2$ on the
raw coefficients. After scaling, the three-block estimator runs in the same
exact rank-one rolling least-squares path as every other linear model in
this paper, refit at every bar. Setting $\beta_3 = 0$ recovers the
two-block model exactly.

Each penalty restates a measured result. $\alpha_1 = 1$:
Section~\ref{sec:dense_weak_twoblock} measured that shrinking the low-dimensional,
fully informative HAR block raises loss, so the backbone is left
essentially unpenalized. $\alpha_2 = 100$: the exogenous signal is dense
and weak, and heavy uniform shrinkage was the measured optimum.
$\alpha_3 = 1000$: the interaction study measured that the product block's
contribution is positive and stable only when it is shrunk an order of
magnitude harder than the linear block; the one-order-of-magnitude gap
between $\alpha_2$ and $\alpha_3$ is the ratio the study's own ladder
selected on its panel. No penalty was chosen by grid search on the
evaluation loss.
